\documentclass[12pt,a4paper]{article}
\usepackage{amsmath,amssymb,amsthm}
\usepackage{hyperref}
\usepackage{geometry}
\geometry{margin=2.5cm}
\usepackage{enumitem}
\usepackage{booktabs}

\newtheorem{theorem}{Theorem}[section]
\newtheorem{lemma}[theorem]{Lemma}
\newtheorem{corollary}[theorem]{Corollary}
\newtheorem{proposition}[theorem]{Proposition}
\theoremstyle{definition}
\newtheorem{definition}[theorem]{Definition}
\theoremstyle{remark}
\newtheorem{remark}[theorem]{Remark}

\title{Baryogenesis from Recognition Science:\\
Sakharov Conditions as Emergent Ledger Properties}

\author{Jonathan Washburn\\
Recognition Science Research Institute\\
Austin, Texas\\
\texttt{washburn.jonathan@gmail.com}}

\date{March 2026}

\begin{document}
\maketitle

\begin{abstract}
We derive baryogenesis---the origin of the matter--antimatter
asymmetry $\eta = n_b/n_\gamma \approx 6.1 \times 10^{-10}$---from
Recognition Science (RS).  Sakharov's three necessary
conditions~\cite{Sakharov1967} emerge structurally from the ledger:
(1)~baryon number violation from defect recombination at early ticks;
(2)~$C$ and $CP$ violation from the double-entry bookkeeping
structure and three-generation mixing
($\mathrm{face\_pairs}(3) = 3$); (3)~departure from thermal
equilibrium from the irreversible dynamics of early ledger epochs.
The baryon asymmetry magnitude $\eta \approx \varphi^{-44}$ is a
$\varphi$-ladder consistency check: rung $-44$ gives
$\varphi^{-44} \approx 6.38 \times 10^{-10}$, compared with the
observed $\eta \approx 6.104 \times 10^{-10}$ (Planck 2018,
$4.6\%$ discrepancy, $\approx 5\sigma$).  The rung $-44$
identification is post-hoc (matched to observation); a
zero-parameter derivation from the RS Hamiltonian is an identified
open problem.  No separate inflaton, grand unified theory, or
leptogenesis mechanism is needed.
\end{abstract}

%% ============================================================
\section{Introduction}
\label{sec:intro}
%% ============================================================

The observed universe contains ${\sim}\,10^{80}$ baryons and
essentially zero antibaryons.  This asymmetry, quantified by the
baryon-to-photon ratio $\eta \approx 6.1 \times 10^{-10}$, must
have been generated dynamically in the early universe---a process
called baryogenesis.

Sakharov~\cite{Sakharov1967} identified three necessary conditions:
\begin{enumerate}
  \item Baryon number ($B$) violation.
  \item $C$ (charge conjugation) and $CP$ (charge-parity) violation.
  \item Departure from thermal equilibrium.
\end{enumerate}

Several mechanisms have been proposed (GUT baryogenesis,
electroweak baryogenesis, leptogenesis, Affleck--Dine), each
requiring specific new physics.  RS~\cite{RS_Axioms} provides all three conditions
from the ledger structure alone.

%% ============================================================
\section{Recognition Science Framework}
\label{sec:framework}
%% ============================================================

All RS results follow from three axioms~\cite{RS_Axioms}:
\begin{enumerate}[label=(A\arabic*),nosep]
  \item $J(1) = 0$ \hfill\textit{(Normalization)}
  \item $J(xy) + J(x/y) = 2J(x)J(y) + 2J(x) + 2J(y)$ for all $x, y > 0$
  \hfill\textit{(Recognition Composition Law)}
  \item $J''_{\log}(0) = 1$, where $J_{\log}(t) := J(e^t)$
  \hfill\textit{(Calibration)}
\end{enumerate}

\begin{theorem}[Cost Uniqueness~\cite{RS_Axioms}]
\label{thm:J}
The unique continuous solution is $J(x) = \tfrac{1}{2}(x + x^{-1}) - 1$,
satisfying $J(x) \geq 0$ with equality iff $x = 1$.
\end{theorem}

\begin{proof}[Proof sketch]
Setting $H(t) = 1 + J(e^t)$ transforms (A2) into the d'Alembert equation
$H(s+t)+H(s-t) = 2H(s)H(t)$ with $H(0)=1$ (A1) and $H''(0)=1$ (A3).
By the Acz\'el classification~\cite{Aczel1966}, $H(t) = \cosh t$, giving
$J(x) = \tfrac{1}{2}(x+x^{-1})-1$.
\end{proof}

Note that $J$ satisfies the \emph{reciprocal symmetry}
$J(x) = J(x^{-1})$ for all $x > 0$.  This symmetry reflects
CPT invariance at the level of the cost functional.  C violation
and CP violation arise not from $J$ itself but from the
\emph{asymmetric phase structure} of the 8-tick epoch and the
asymmetric initial conditions (high defect at $t = 0$), as
explained in Sections~\ref{sec:c-violation}--\ref{sec:cp-violation}.

%% ============================================================
\section{Sakharov Condition 1: $B$ Violation}
\label{sec:b-violation}
%% ============================================================

\begin{proposition}[Baryon Number Violation --- Structural Argument]
\label{thm:b-violation}
The ledger topology permits $B$-changing processes through defect
recombination at early ticks.
\end{proposition}

\begin{proof}
In the early ledger (few epochs after $t = 0$), the defect density
is high and the voxel configuration is far from equilibrium.
Recognition events can recombine in topologically non-trivial ways:
three defects (one per face pair of the $D = 3$ cube) can merge
into a single lower-energy configuration, changing the net baryon
number by $\Delta B = 1$.  This is the ledger analog of sphaleron
processes in the Standard Model, where the $SU(2)_L$ instanton
mediates $\Delta B = \Delta L = 3$ transitions.
\end{proof}

\begin{remark}
At late times (after the electroweak phase transition), the defect
density drops and the recombination rate is exponentially suppressed:
$\Gamma_B \propto e^{-E_{\mathrm{def}}/T}$ where $E_{\mathrm{def}}$
is the defect energy.  Baryon number is effectively conserved in the
current epoch, consistent with proton stability observations
($\tau_p > 10^{34}\;\mathrm{yr}$).
\end{remark}

%% ============================================================
\section{Sakharov Condition 2: $C$ and $CP$ Violation}
\label{sec:c-violation}
\label{sec:cp-violation}
%% ============================================================

\begin{remark}[Reconciling $C$ Violation with $J(x) = J(x^{-1})$]
\label{rem:c-reconcile}
The reciprocal symmetry $J(x) = J(x^{-1})$ of the cost functional
corresponds to invariance under particle-antiparticle exchange at the
\emph{static} cost level.  This is consistent with CPT invariance
(proved from $J(x)=J(x^{-1})$ in the companion paper on CP
violation).  However, $C$ symmetry can still be violated dynamically
by:
\begin{enumerate}[nosep,label=(\roman*)]
  \item The \emph{asymmetric phase ordering} of the 8-tick epoch:
  the 8 phases $(0,1,\ldots,7)$ are ordered, and $C$ acts on these
  phases in an orientation-dependent way that depends on the
  specific 3-bit decomposition.
  \item The \emph{asymmetric initial conditions}: the high-defect
  early ledger ($t \approx 0$) breaks time-reversal symmetry, and
  by the CPT theorem, breaking T implies breaking CP (since CPT is
  preserved).
  \item The \emph{three-generation CKM phase}: with $N = 3$
  generations, there is an irreducible $CP$-violating phase
  $\delta_{\mathrm{CKM}}$ in the mixing matrix.
\end{enumerate}
The $J$-cost symmetry $J(x) = J(x^{-1})$ guarantees CPT is preserved,
but it does not guarantee C or CP separately.
\end{remark}

\begin{proposition}[$C$ Violation from Double-Entry Bookkeeping --- Structural Identification]
\label{thm:c-violation}
The ledger's double-entry bookkeeping distinguishes particles from
antiparticles, providing a structural basis for $C$ violation.
\end{proposition}

\begin{proof}[Structural argument]
Each recognition event records a debit-credit pair on the ledger.
Charge conjugation $C$ maps $x_e \mapsto x_e^{-1}$ on individual
ledger entries (particle $\to$ antiparticle), which preserves $J$
by the reciprocal symmetry.  However, the \emph{phase ordering}
within the 8-tick epoch assigns tick indices $k = 0, 1, \ldots, 7$
to successive phases.  Charge conjugation combined with time
reversal (CT) would map $k \mapsto 8 - k$, which is not the same
as the original ordering.  Since CPT is preserved but T is broken
(early-universe monotone defect decrease), C must also be broken,
as $(\text{CPT}) = (\text{CP}) \times \text{T}$ implies
$(\text{CP}) = (\text{CPT})/\text{T}$.  With T broken, CP and C
each break at rates governed by the CKM phase $\delta$.
\end{proof}

\begin{theorem}[$CP$ Violation from Three Generations]
\label{thm:cp}
With $N_{\mathrm{gen}} = 3$ generations (from
$\mathrm{face\_pairs}(3) = 3$), the CKM mixing matrix has an
irreducible $CP$-violating phase.
\end{theorem}

\begin{proof}
A unitary $N \times N$ mixing matrix has
$(N-1)(N-2)/2$ independent $CP$-violating phases.  For $N = 3$:
$(3-1)(3-2)/2 = 1$ phase.  For $N = 2$: zero phases (no $CP$
violation).  The third generation is essential for $CP$ violation.
Since $N_{\mathrm{gen}} = 3$ is forced by the cube structure,
$CP$ violation is a structural consequence.
\end{proof}

\begin{remark}
The Jarlskog invariant $J_{\mathrm{CP}} \approx 3.0 \times 10^{-5}$
quantifies the physical $CP$ violation.  In RS, this is derived from
the $\varphi$-ladder structure of the CKM matrix (see the companion
paper on the CKM matrix).
\end{remark}

%% ============================================================
\section{Sakharov Condition 3: Out of Equilibrium}
\label{sec:out-of-eq}
%% ============================================================

\begin{proposition}[Departure from Equilibrium --- Structural Argument]
\label{thm:out-eq}
The early ledger dynamics are intrinsically out of equilibrium in
the sense relevant for baryogenesis.
\end{proposition}

\begin{proof}
At the earliest epochs, the ledger has many entries far from unity,
corresponding to high total defect $D \gg 0$.  (The zero-defect
state $\mathbf{1}$ is the asymptotic \emph{future} equilibrium of RS
dynamics, not the initial state---see the companion paper on the
arrow of time.)  This high-defect early state is far from
thermodynamic equilibrium: defect is decreasing rapidly through
irreversible recognition events, and the rate of defect reduction
(analogous to the Hubble expansion rate) exceeds the rate at which
$B$-violating reactions can equilibrate.  During this period,
$B$-violating and $CP$-violating processes operate at rates different
from their inverses, generating a net baryon asymmetry.
``Equilibrium'' here means detailed balance of reaction rates, not
minimum cost; the early universe lacks such balance because the rapid
descent of defect prevents the forward and reverse $B$-changing rates
from equilibrating.
\end{proof}

%% ============================================================
\section{Baryon Asymmetry Magnitude}
\label{sec:magnitude}
%% ============================================================

\begin{proposition}[$\varphi$-Ladder Consistency Check for $\eta$]
\label{thm:eta}
The observed baryon-to-photon ratio $\eta \approx 6.1 \times 10^{-10}$
falls near rung $-44$ of the $\varphi$-ladder:
\begin{equation}
  \boxed{\eta \approx \varphi^{-44} \approx 6.38 \times 10^{-10}.}
\end{equation}
\end{proposition}

\begin{proof}
Rung $-44$: $\varphi^{-44} = (1/\varphi)^{44} \approx (0.6180)^{44}$.
Using $\varphi^{10} \approx 122.99$:
$\varphi^{44} = \varphi^{40} \times \varphi^{4} \approx (122.99)^4 \times 6.854
\approx 2.285 \times 10^8 \times 6.854 \approx 1.566 \times 10^9$.
Therefore $\varphi^{-44} \approx 6.38 \times 10^{-10}$.
\end{proof}

\begin{remark}[Status of the $\eta$ prediction]
\label{rem:eta-status}
Rung $-44$ is the $\varphi$-ladder value closest to the observed
$\eta \approx 6.1 \times 10^{-10}$.  The identification $\eta \approx
\varphi^{-44}$ is a consistency check, not a zero-parameter derivation:
the rung number $44$ is obtained by matching to observation, not by
an independent RS calculation of the Jarlskog invariant
$J_{\mathrm{CP}}$, the out-of-equilibrium factor, and the $B$-violation
efficiency from first principles.  A genuine zero-parameter RS
derivation of $\eta$ would require computing these three factors from
the ledger Hamiltonian without reference to the observed value; this
is an identified open problem.
\end{remark}

\begin{remark}[Comparison with Planck]
The Planck measurement:
$\eta = (6.104 \pm 0.058) \times 10^{-10}$~\cite{Planck2020}.
The RS $\varphi$-ladder consistency value $\varphi^{-44} \approx
6.38 \times 10^{-10}$ agrees to within
$(6.38 - 6.10)/6.10 \approx 4.6\%$, or approximately $5\sigma$
given the Planck precision.  While this is a reasonable order-of-magnitude
match given the post-hoc nature of the rung identification, it is
not a precise zero-parameter prediction.
\end{remark}

%% ============================================================
\section{Comparison with Standard Mechanisms}
\label{sec:comparison}
%% ============================================================

\begin{center}
\renewcommand{\arraystretch}{1.3}
\begin{tabular}{@{}lccc@{}}
\toprule
\textbf{Mechanism} & \textbf{New physics} &
\textbf{$\eta$ predicted?} & \textbf{Status} \\
\midrule
GUT baryogenesis & Superheavy bosons & Qualitative & Untested \\
EW baryogenesis & 1st-order EWPT & Insufficient $CP$ & Marginal \\
Leptogenesis & Heavy $\nu_R$ & Qualitative & Untested \\
Affleck--Dine & Flat directions & Qualitative & SUSY-dependent \\
RS & None & $\varphi^{-44}$ (consistency) & $4.6\%$ off Planck \\
\bottomrule
\end{tabular}
\end{center}

%% ============================================================
\section{Predictions}
\label{sec:predictions}
%% ============================================================

\begin{enumerate}
  \item \textbf{$\eta \approx \varphi^{-44}$}: A $\varphi$-ladder
  consistency check placing $\eta$ at rung $-44$.  A fully
  independent RS derivation requires computing the Jarlskog invariant,
  out-of-equilibrium factor, and $B$-violation efficiency from the
  ledger Hamiltonian; this derivation is an identified open problem.

  \item \textbf{No neutron-antineutron oscillation}: Since baryon
  number violation is exponentially suppressed in the current epoch,
  $n$-$\bar{n}$ oscillation should not be observed.

  \item \textbf{No antihelium in cosmic rays}: The baryon asymmetry
  is global; no primordial antimatter domains exist.
\end{enumerate}

%% ============================================================
\section{Conclusion}
\label{sec:conclusion}
%% ============================================================

Baryogenesis in Recognition Science is a structural consequence of
the ledger: all three Sakharov conditions emerge from the discrete
ledger structure without new physics.  The $\varphi$-ladder assigns
rung $-44$ to the baryon-to-photon ratio, giving
$\varphi^{-44} \approx 6.38 \times 10^{-10}$, consistent with
Planck's $\eta = 6.104 \times 10^{-10}$ to $4.6\%$.  A zero-parameter
derivation of $\eta$ from the RS Hamiltonian (computing the Jarlskog
invariant, out-of-equilibrium factor, and $B$-violation efficiency
from first principles) is an identified open problem.

%% ============================================================
\begin{thebibliography}{99}

\bibitem{Aczel1966}
J.~Acz\'el, \textit{Lectures on Functional Equations and Their
Applications}, Academic Press, New York (1966).

\bibitem{RS_Axioms}
J.~Washburn, ``The Algebra of Reality: A Recognition Science Derivation
of Physical Law,'' \emph{Axioms} \textbf{15}(2), 90 (2026).
\texttt{doi:10.3390/axioms15020090}

\bibitem{Sakharov1967}
A.~D.~Sakharov, ``Violation of CP Invariance, C Asymmetry, and
Baryon Asymmetry of the Universe,'' JETP Lett.\ \textbf{5}, 24
(1967).

\bibitem{Planck2020}
Planck Collaboration, ``Planck 2018 Results.\ VI.\ Cosmological
Parameters,'' Astron.\ Astrophys.\ \textbf{641}, A6 (2020).

\end{thebibliography}

\end{document}
